\section{Zoo DAO}

At the heart of the Zoo Protocol lies the Zoo \DAO{}, a community-driven initiative dedicated to preserving endangered species from the brink of extinction. The \DAO{} offers rewarding experiences and empowers members to shape the future of the Zoo Protocol through meaningful decisions.

Zoo \DAO{} governs the Zoo Protocol, a sophisticated set of smart contracts for collateral-backed NFT Animal assets. When users join the Zoo \DAO{}, they earn Experience Points (XP) and governance tokens (KEEPER), empowering them to participate actively in the protocol.

\subsection{Experience Points and Leveling}

The accrual of XP tokens is tied to undertaking ``quests'' or tasks, each with a specific reward allocation of XP and KEEPER tokens. These quests vary in complexity, with higher-level tasks requiring more substantial contributions to the community and yielding more XP. The more XP earned, the faster one advances through the ranks. Level is determined by the equation:
\begin{equation}
\text{Level} = \text{level} \times 300 - 500
\end{equation}

Progressing through levels grants greater authority within the ecosystem and broadens influence. XP tokens are non-transferable and strictly tied to Decentralized Identification Services (DIS) and comprehensive KYC procedures, preserving the integrity of the system.

\subsection{Governance and Voting}

Votes are open for all members; however, some voting privileges are reserved for members of certain levels. Generally, higher-level votes are more privileged and controlling decisions. Specific actions such as proposing, voting, and interacting with governance mechanisms may require a blend of \$ZOO and KEEPER tokens. Different proposals cater to different levels, underlining the ecosystem's commitment to inclusivity.

\subsection{Betting-Based Governance}

Central to the Zoo Governance and Voting Protocol is the vZOO level and the stake age duration. Unique to the system, KEEPER token holders ``bet'' on proposals rather than voting. Winning bets yield additional tokens, promoting active participation. The process allows anyone who meets the requirements to submit a proposal, with those gaining traction advancing to the next round, while lesser-supported ones are eliminated. This distinctive approach ensures a vibrant and thriving ecosystem, rewarding active participation and fostering continuous community involvement.
